\chapter{Thesis Materials and Sources}
\label{chap:tmas} Many surface ozone (O$_3$) estimations have been developed using air pollution measurements, meteorological imagery, and high-resolution estimations through various combinations of these (e.g. \cite{alvarez-mendoza_spatial_2019}, \cite{schultz_tropospheric_2017}, \cite{wang_machine_2022}). Methods such as kriging and complex trend estimation at high temporal resolutions (i.e., hourly, daily, and weekly) are constrained by the availability of imagery, which is influenced by orbital patterns and geographical characteristics of the Earth. The associated imagery and additional band information presented here illustrate the extensive range of imagery now available through Big Data. The National Oceanic and Atmospheric Administration/National Centers for Environmental Prediction (NOAA/NCEP), the European Centre for Medium-Range Weather Forecasts (ECMWF) and the National Aeronautics and Space Administration/Global Modeling and Assimilation Office (NASA/GMAO) provide rich ensembles of climate data products, enabling comprehensive daily coverage within this thesis (\cite{schultz_tropospheric_2017}). These datasets can be further combined in overlapping areas due to shared interests among nations (\cite{gaudel_tropospheric_2018}; \cite{carrillo-torres_use_2017}; \cite{xu_effectiveness_2022}). 

This section provides details on datasets used after filtering through myriads of features and other materials available for future work. The features used in this project were selected from many initial features collected before filtering the training values to the final time period. The statistical distribution of these are in Table \ref{tab:variable_codes}. The standard deviation (S.Dev) for each variable indicates the overall width of the feature compared to the overall ranges of features gathered in Chapter \ref{methods:varextraction}. \begin{table}[h!]
    \centering
    \begin{tabular}{ccccccc}
        \hline
        \textbf{$\beta_{n}$} & \textbf{Variable Name} & \textbf{Min} & \textbf{Median} & \textbf{Max} & \textbf{Mean}  & \textbf{S.Dev} \\
        \hline
        $y(\beta)$ & DAMO$_3$ & 0.048	& 0.012 & 0.004 & 0.108 & 0.048 \\
        $\beta_{1}$ & Specific Humidity & 0.006	& 0.003 & 0 & 0.018 & 0.005 \\
        $\beta_{2}$ & D.S Radiation & 237.26	& 81.172 & 37.991 & 375.957 & 240.381 \\
        $\beta_{3}$ & Min Surface Temp. & 288.136	& 8.56 & 265.733 & 307.929 & 287.645 \\
        $\beta_{4}$ & Max Surface Temp. & 303.133	& 9.016 & 274.183 & 321.976 & 303.707 \\
        $\beta_{5}$ & Burn Index & 52.498	& 26.95 & -479.182 & 148.454 & 56.055 \\
        $\beta_{6}$ & Mean Vapor Deficit & 2.481	& 1.425 & 0 & 6.846 & 2.239 \\
        $\beta_{7}$ & NDVI & 0.224	& 0.085 & 0.092 & 0.66 & 0.202 \\
        $\beta_{8}$ & Stratospheric NO$_2$ & 0 & 0 & 0 & 0 & 0 \\
        $\beta_{9}$ & S5P 1km TCO$_3$ & 0.133	& 0.012 & 0.099 & 0.208 & 0.131 \\
        $\beta_{10}$ & S5P TCO$_3$ Temperature & 229.042	& 3.022 & 219.783 & 235.292 & 229.567 \\
        $\beta_{11}$ & Carbon Monoxide & 0.028	& 0.006 & -0.139 & 0.228 & 0.028 \\
        $\beta_{12}$ & Water Column Density & 999.735	& 775.724 & -7,240.62 & 15,846.62 & 728.739 \\
        $\beta_{13}$ & Water Column Energy & 4.604	& 3.652 & -32.304 & 73.623 & 3.305 \\
        $\beta_{14}$ & Formaldehyde () & 0 & 0	& -0.002 & 0.021 & 0 \\
        $\beta_{15}$ & CH$_2$O Slant Column & 0 & 0	& -0.003 & 0.015 & 0 \\
        $\beta_{16}$ & Cloud Fraction & 0.19	& 0.253 & -0.276 & 1.073 & 0.082 \\
        $\beta_{17}$ & Dobson KE & 380.55	& 149.055 & 49.36 & 754.329 & 375.743 \\
        $\beta_{18}$ & S5P KE & 599.759	& 228.672 & 80.41 & 1,215.22 & 589.601 \\
        $\beta_{19}$ & $\beta_{2}$ WkAvg & 237.265	& 76.921 & 71.287 & 365.422 & 241.321 \\
        $\beta_{20}$ & Wind Velocity WkAvg & 2.783	& 0.819 & 0.624 & 7.319 & 2.762 \\
        $\beta_{21}$ & $\beta_{6}$ WkAvg & 2.48	& 1.347 & 0.117 & 6.443 & 2.298 \\
        $\beta_{22}$ & Dobson WkAvg & 0.063	& 0.005 & 0.05 & 0.087 & 0.063 \\
        $\beta_{23}$ & $\beta_{3}$ WkAvg & 303.124	& 8.628 & 278.99 & 320.202 & 303.666 \\
        $\beta_{24}$ & $\beta_{11}$ WkAvg & 229.039	& 2.911 & 221.381 & 233.903 & 229.556 \\
        $\beta_{25}$ & Spring & 0 & 0 & 1 & SEA & SON \\
        $\beta_{26}$ & Summer & 0 & 0 & 1 & SEA & SON \\
        $\beta_{27}$ & Winter & 0 & 0 & 1 & SEA & SON \\
       \hline
    \end{tabular}
    \caption{Feature codes and Numerical Distributions.}
    \label{tab:variable_codes}
\end{table} The connotation WkAvg denotes weekly moving averages of the associated $\beta_{n}$ or metric from which it was derived. The databases from which these features were sourced contain several additional bands that can be leveraged for future pollution models. As discussed in this chapter, these are encapsulated within a cloud-based Big Data source hosted by Google, known as Earth Engine (GEE). This engine enables the integration of RK methodology through JavaScript development into contemporary GIS ML/AI modeling, providing powerful solutions for Big Data and future projects (\cite{xu_mapping_2023}). Updates to the current coding libraries in Python and JavaScript were essential for ensuring the reproducibility of this project and facilitated the implementation of methods discussed in Chapter \ref{methods:inrtp} through Python. 
%%%%%%%%%%%%%%%%%%%%%%%%%%%%%%%%%%%%%%%%%%%%%%%%%%%%%%%%%%%%%%%%%%%%%%%%%%%%%%%%%%%%%%%%%%%%%%%%%%%%%%%%%%%%%%%%%%%%
\section{Vector Data} \label{tmas:vdata} Vector data in GIScience comprises points, lines, and polygons. This thesis utilizes points for monitoring locations and polygons for the census tracts/overall shape of PHOTUC. In Python, EPA monitor data was processed to yield daily maximum 8-hour surface O$_3$ (DAMO$_3$) values, along with the corresponding location and date, in alignment with the literature that typically adopts this metric (e.g., \cite{flynn_spatial_2021}; \cite{runmei_ma_full-coverage_2021}; \cite{sofen_gridded_2015}). The distribution of these can be seen in Figure \ref{fig:monsurfo3hist} where the overall distribution is normal. To assess potential inconsistent trends at both ends of the dataset, these values are divided into the 5th and 95th percentiles. \begin{figure}[h!]  
    \centering
    \includegraphics[width=0.8\linewidth]{ThesisFilePkg//figs/surf_o3_hist_by_site_2019_2024.png}
    \caption{Histogram of Daily Average Maximum 8-Hour O$_3$ (DAMO$_3$) at Monitor Locations.}
    \label{fig:monsurfo3hist}
\end{figure} Geographic and temporal information was used to extract data from the gathered images. Dummy temporal variables were created through one-hot encoding to impose temporal restrictions on each model. Month and seasonal values were inspected, reflecting established temporal patterns identified in the existing literature (\cite{shan_observational_2008}; \cite{abdullah_development_2019}; \cite{balamurugan_tropospheric_2021}). While the date served as a filtering feature for the Residual Kriging method, it was not employed during the training and testing of the statistical models within the complex ensemble training phase. Incorporating the date as a feature would have resulted in nearly 365 dummy variables or improper grouping variables, thereby increasing computation time and introducing a bias towards yearly indexing in the models. This topic is further explored in Chapter \ref{rlts:modelacc}.  
%%%%%%%%%%%%%%%%%%%%%%%%%%%%%%%%%%%%%%%%%%%%%%%%%%%%%%%%%%%%%%%%%%%%%%%%%%%%%%%%%%%%%%%%%%%%%%%%%%%%%%%%%%%%%%%%%%%%
\subsection{Arizona Boundaries} \label{tmas:azbounds} The Arizona state boundary was downloaded from the TIGER line Census (\cite{united_states_census_bureau_tigerline_2022}). The shapefile and related informational files are a subset of select geographic information from the U.S. Census Bureau's Master Address File. These boundaries are depicted in Figure \ref{fig:basemap} created prior to the ensemble training phase. \begin{figure}[h!]
    \centering
    \includegraphics[width=0.95\linewidth]{ThesisFilePkg//figs/maps/BaseAOI.png}
    \caption{Monitor Locations Within the Area of Interest (AOI) in Arizona}
    \label{fig:basemap} 
\end{figure} A section within this vase database contains the Topologically Integrated Geographic Encoding and Referencing (TIGER) Database (MTDB). The MTDB represents a seamless national file with no overlaps or gaps between parts. This feature makes it an ideal choice for projects, as each data product can be combined to create a comprehensive dataset for nationwide analysis. 

The boundaries of incorporated places for this time period in the Arizona shapefile have been current since January 1, 2021, as reported by the US Census Bureau's Boundary and Annexation Survey (BAS) (\cite{united_states_census_bureau_arizona_2025}). Shapefile data were downloaded, filtered, and displayed at both the state and county levels (Figure \ref{fig:basemap}) to establish the general limits for monitor selection within the area of interest (AOI). Each state and county-level boundary file underwent a unification process to eliminate inner state boundaries, resulting in a singular representation of the study region that includes Phoenix and Tucson (PHOTUC). The boundary consisting of the counties of Pinal, Pima, and Maricopa was utilized as a reference to incorporate both cities. 
%%%%%%%%%%%%%%%%%%%%%%%%%%%%%%%%%%%%%%%%%%%%%%%%%%%%%%%%%
\subsection{EPA O$_3$ Monitors} \label{tmas:mondata} The EPA provides data on various ground-level pollutants and their associated meteorological drivers (\cite{epa_air_2025}), offering valuable opportunities discussed later in Chapter \ref{chap:clcsn}. Initially, daily monitoring data were collected from January 1, 2018, to the present day. The initial final date is noted for future work, allowing for potential stopping dates up to the present day. The daily average maximum O$_3$ (DAMO$_3$) was utilized, as similar studies investigating O$_3$ concentrations have shown that it better reflects the health impacts of exposure and transport (e.g., \cite{adebayo-ojo_short-term_2022}, \cite{carbajal-arroyo_effect_2011}, \cite{fine_variability_2015}, \cite{maji_prescribed_2024}). Geospatial information contained in the monitoring metadata was used to filter all available data points in Arizona during this period. The distribution of features at each monitor used in this thesis is illustrated in Figure \ref{fig:histograms}. \begin{figure}[h!]
    \centering
    \includegraphics[width=0.75\linewidth]{ThesisFilePkg//figs/hist_all_2019_2024.png}
    \caption{Potential Features and Distributions at Monitoring Locations}
    \label{fig:histograms}
\end{figure} Each monitor's location and DAMO$_3$ mixing ratio were extracted and averaged for each available day over a five-year period. This resulted in 2,215 potential daily averages per monitoring site. This API also offers access to a wide range of air quality data, including daily and hourly measurements of various pollutants other than O$_3$. A total of 65 monitors in Arizona were identified through unique iterations of the provided site number from 1980 to the present day. In 2023, 38 of these monitors were located within Maricopa, Pinal, and Pima counties \cite{epa_air_2025}. Only model perfect monitors were utilized, leaving the remaining monitors for AOI expansion in later work.

The distribution of monitors is illustrated in Figure \ref{fig:basemap} within the AOI. The dates associated with each monitor were used for raster extraction, with Google Earth Engine (GEE) enabling the rapid assignment of constraint variables to these dates (\cite{cardille2024}). In the PHOTUC region, there were 32 non-moved monitors, of which 28 were complete monitors over the 5-year period. Further development of interpolation and imputation strategies, mentioned in Chapter \ref{dev} can better account for these missing monitors and subsequent features. Given the temporal limitations of orbital patterns, some features  utilized the smoothed weekly series described in Chapter \ref{methods:miar}. As a result, some incomplete weeks near the ends of the dataset were removed for data such as the normalized difference vegetation index (NDVI). A total of 60,900 samples are utilized in each model from January 1st, 2019 to December 14th, 2024.
%%%%%%%%%%%%%%%%%%%%%%%%%%%%%%%%%%%%%%%%%%%%%%%%%%%%%%%%%%%%%%%%%%%%%%%%%%%%%%%%%%%%%%%%%%%%%%%%%%%%%%%%%%%%%%%%%%%%
\section{Raster Data} \label{tmas:satdat} Most of this project utilized raster data, an exceptionally well-structured mathematical representation of signals. Otherwise known as raster type data in GIScience, it serves as the primary output in remote sensing, a field focused on Earth-bound, sun-sourced light captured by satellites (\cite{jensen2007}. These representations, commonly referred to as "rasters" in geography, consist of a grid of values depicted as pixels, which can represent various phenomena. When these values originate from a signal, the mathematical transformations discussed in Chapter \ref{methods:ppro} can be applied to accurately represent their metaphorical value( \cite{jensen2007}). 

For instance, temperature values can be estimated with extreme precision based on the differences in near-infrared (NIR) and incoming UV radiation. Much of this data is available in GEE's extensive cloud database. GEE was accessed via Python to establish the framework for the images utilized in the final predictive model. GEE is designed for planetary-scale environmental data analysis (\cite{gharibvand_changes_2023}) and offers access to a comprehensive repository of satellite imagery and geospatial datasets, along with Java-based tools for analysis and visualization (\cite{gorelick_google_2017}). 
%%%%%%%%%%%%%%%%%%%%%%%%%%%%%%%%%%%%%%%%%%%%%%%%%%%%%%%%%
\subsection{Google Earth Engine Image Collection Exportation} \label{tmas:geeexport} While the interactive code editor and Graphical User Interface (GUI) make it accessible to users proficient in JavaScript, it may pose challenges for those less familiar with coding, such as researchers. However, GEE enables collaborative work by allowing users to easily share scripts, data, and results with those willing to learn. The platform has also established robust Python integration, developed by Dr. Wu at the University of Tennessee (\cite{wu_geemap_2020}). This integration has streamlined the automation process, enabling the collection of scientifically precise images at a resolution of 500 m, with baseline EPSG 4326 metadata through typical remote sensing transformations (\cite{gorelick_google_2017}). Thanks to numerous open-source documents provided by Dr. Wu and his associates, there is no longer a need to learn an additional programming language for geo-statistical modeling, as these resources effectively translate GEE functions from JavaScript into Python. \begin{figure}[h!]
    \centering
    \includegraphics[width=0.85\linewidth]{ThesisFilePkg//figs/geendvi.png}
    \caption{Google Earth Engine (GEE) Graphical User Interface (GUI) In JavaScript.}
    \label{fig:gee}
\end{figure} Figure \ref{fig:gee} shows the code editor associated with GEE's GUI, along with a representation of the normalized difference vegetation index (NDVI) using Landsat 9. GEE facilitates near-real-time (NRT) data analysis in collaboration with institutions such as NASA, which is crucial for monitoring and responding to environmental events as they occur. When accurately translated from Python to JavaScript, SMaRK can operate directly within GEE’s code editor, providing an essential NRT DAMO$_3$ model for the platform. 

GEE provides free access to a vast archive of satellite imagery and geospatial datasets, including Landsat, Sentinel, MODIS, and more. This accessibility is available to users in both Python and R. GEE has gained prominence in recent environmental monitoring, research, and management due to its user-friendly interface. It serves as a key resource for historical analyses, and integration with JavaScript is a goal outlined in Chapter \ref{chap:clcsn}. The features over time are illustrated in Figure \ref{fig:scatters} and were used to assess the quality of data discussed in this chapter. \begin{figure}[h!]
    \centering
    \includegraphics[width=0.85\linewidth]{ThesisFilePkg//figs/scatter_2018_2023.png}
    \caption{Linear Correlation of Features with Surface O$_3$}
    \label{fig:scatters}
\end{figure} The GEE cloud-based infrastructure facilitates large-scale data processing without the need for local high-performance computing resources. This capability is particularly advantageous for the rapid and efficient processing and analysis of extensive data sets.
%%%%%%%%%%%%%%%%%%%%%%%%%%%%%%%%%%%%%%%%%%%%%%%%%%%%%%%%%%%%%%%%%%%%%%%%%%%%%%%%%%%%%%%%%%%%%%%%%%%%%%%%%%%%%%%%%%%
\section{European Re-Analysis: Fifth-Generation} \label{tmas:era5} The fifth generation of European Re-Analysis (ERA5) is a climate reanalysis dataset produced by the European Center for Medium-Range Weather Forecasts (ECMWF) as part of the Copernicus Climate Change Service (C3S) \parencite{hoffmann2019}. ERA5 is one of the latest climate reanalyses produced by ECMWF, providing hourly data on various atmospheric, land-surface, and sea-state parameters, along with estimates of uncertainty. The data is readily accessible from NCAR, featuring grids at a 1 km resolution and atmospheric parameters on 37 pressure levels \parencite{hersbach2020}. The temporal estimates of the dataset extend from 1950 to the present day, with daily updates available 5 days behind real time. ERA5 offers a consistent view of the evolution of land variables over several decades, providing enhanced resolution compared to other datasets, such as LMCS and GRID-Met, which are relevant for meteorological data \parencite{hersbach2020,munoz-sabater2021,tarek2020}.

The asset utilized in this project is a daily aggregate of modeled data combined with satellite observations, creating a globally comprehensive and consistent data set based on the laws of physics for future work. Given that this available datasets spans several decades, it has demonstrated accurate descriptions of climate change over time. This data set includes 50 meteorological variables available from 1950 to three months from today \parencite{european_centre_for_medium-range_weather_forecasts_era5_2019}. These data have been available at a coarse 11 km resolution since the assimilation of satellite technology. The resulting rasters were aggregated to a resolution of 500 m before being exported from GEE. This project utilized the 2 m surface temperature, Earth's surface reflectivity, shortwave radiation at the Earth's surface, net change in solar radiation, northeast/eastern components of wind up to a height of 10 m, surface pressure, and total precipitation accumulation, as these were all suitable precursors for O$_3$ reactions \parencite{zhang_interaction_2021}. 
%%%%%%%%%%%%%%%%%%%%%%%%%%%%%%%%%%%%%%%%%%%%%%%%%%%%%%%%%%%%%%%%%%%%%%%%%%%%%%%%%%%%%%%%%%%%%%%%%%%%%%%%%%%%%%%%%%%
\section{TOMS/OMI Merged 100 KM Surface O$_3$} \label{tmas:era5} The Total Ozone Mapping Spectrometer (TOMS) data is produced and maintained by the Laboratory for Atmospheres at NASA's Goddard Space Flight Center. It yields a long-term, continuous record of the United States’ role in monitoring global and regional trends of surface ozone values over the past twenty-five years. The Ozone Monitoring Instrument (OMI), launched aboard the Aura satellite in July 2004, offers 1.0$\deg$ resolution and was used to calibrate the data collected from the TOMS instrument. Together, these instruments create a merged ozone product, integrating observations from TOMS/EarthProbe, TOMS/Nimbus-7, TOMS/Meteor-3, and OMI/Aura, along with USGS-based interpolations for periods where data is lacking. This is the coarsest dataset in this thesis, and the aggregation to 500 m from GEE is likely to reveal very common values in a gradient-like structure. This dataset employs transformations from Chapter \ref{methods:tomi} to enable better trends with ERA5 data in acquisition periods from 1978 to the present. The TOMS/OMI \ang{1.0} resolution, when converted to meters, is approximately 111 KM.
%%%%%%%%%%%%%%%%%%%%%%%%%%%%%%%%%%%%%%%%%%%%%%%%%%%%%%%%%%%%%%%%%%%%%%%%%%%%%%%%%%%%%%%%%%%%%%%%%%%%%%%%%%%%%%%%%%%%
\section{Sentinel-5P/TROPOMI Products} \label{tmas:S5P} In 2017, CAMS launched the Sentinel-5 Precursor (S5P), the first Copernicus-based mission dedicated to monitoring the planet’s atmosphere \parencite{veefkind2012}. Addressing the need for improved data to analyze the environmental concerns discussed in Chapter \ref{introchap:whyO$_3$}, the Copernicus Atmosphere Monitoring Service (CAMS) has emerged as a leading institution for monitoring greenhouse gases. Its primary objective is to provide policymakers, businesses, and citizens with enhanced atmospheric environmental information \parencite{petetin_model_2021}. The European Centre for Medium-Range Weather Forecasts (ECMWF), an independent intergovernmental organization supported by 35 states, implements CAMS on behalf of the European Union for these purposes.

The sensor on S5P used in this thesis is the TROPOspheric Monitoring Instrument (TROPOMI). It contains a vast array of pollutant archives from it's launch midway through 2017. For most of the allotted pollutants, there are two products available: Near Real-Time (NRT) and Offline (OFFL). The NRT asset covers a smaller area but are made available daily. The OFFL assets contain data from a single orbit and have a 3 day lag, but contains data for a single hemisphere. This leads to necessary interpolation strategies used in Chapter \ref{methods:inrtp} where data was missing at everyother interval due to these development schemes. TROPOMI products are the primary source of tropospheric O$_3$ monitoring today \parencite{keppens2024} and offer a vast array of high quality raster products. 

These products have undergone multiple validations and are extensively used in prominent studies at both global and national levels due to their relative high precision, extensive geographic coverage, and long time frame \parencite{zhang_ozone_2019}. This instrument offers a multitude of atmospheric variables including aerosol, NO$_x$, and O$_3$ measurements which are used in this project. Only VOCs/pollutants that are recognized as significant drivers of O$_3$ and spatially relevant were selected from the S5P satellite. Precursor aerosols, such as SO$_2$ concentrations, tend to be localized, making them difficult to acquire in GEE for daily temporal resolutions.
%%%%%%%%%%%%%%%%%%%%%%%%%%%%%%%%%%%%%%%%%%%%%%%%%%%%%%%%%
\subsection{Aerosol Index}  \label{tmas:aai} This data set from the TROPOMI instrument provides high-resolution imagery of the UV Aerosol Index (UVAI) or the Absorbing Aerosol Index (AAI). The AAI is determined based on wavelength-dependent changes in Rayleigh scattering observed in the UV spectrum for a specific pair of wavelengths. It is calculated as the difference between observed and modeled reflectance \parencite{degraaf2005}. The wavelengths of 354 nm and 388 nm used in this index have very low ozone (O$_3$) absorption, allowing AAI to be calculated in the presence of clouds, unlike aerosol optical thickness measurements. This ability enables daily global coverage of potential DAMO$_3$ within the index. These rasters are available with daily averages from GEE and the CAMS home website at a resolution of about 1 km. The AAI represents a measure of the prevalence of aerosols in the atmosphere, but was found to be weighted towards other aerosols given its low corelation with DAMO$_3$.
%%%%%%%%%%%%%%%%%%%%%%%%%%%%%%%%%%%%%%%%%%%%%%%%%%%%%%%%%
\subsection{Nitrogen Dioxide (NO$_2$) Retrievals} \label{tmas:no2} Nitrogen oxides (NO$_2$ and NO) are important trace gases in the atmosphere of Earth, present both in the troposphere and the stratosphere. They enter the atmosphere through anthropogenic activities, such as fossil fuel combustion and biomass burning, as well as natural processes like wildfires, lightning, and microbiological activities in soils \parencite{zhang_characteristics_2024}. NO$_2$ values represent the collective concentrations of nitrogen oxides. This measurement accurately reflects the daytime photochemical cycle involving O$_3$. The NO$_2$ cycle converts NO into NO$_2$ and vice versa on a time scale of minutes. The stratospheric vertical column of the NO$_2$ and NO$_2$ slant column density bands provides an estimate of the O$_3$ reactions occurring in the atmosphere.
%%%%%%%%%%%%%%%%%%%%%%%%%%%%%%%%%%%%%%%%%%%%%%%%%%%%%%%%%
\subsection{Carbon Monoxide (CO) and H$_2$O Column Number Densities} \label{tmas:co} Carbon monoxide (CO) is a crucial atmospheric trace gas for understanding tropospheric chemistry. In certain urban areas, it serves as a significant atmospheric pollutant. The primary sources of CO include the combustion of fossil fuels, biomass burning, and the atmospheric oxidation of methane and other hydrocarbons. While fossil fuel combustion is the predominant source of CO in the northern mid-latitudes, the oxidation of isoprene and biomass burning play a vital role in the tropics. TROPOMI on the S5P satellite monitors the global abundance of CO by utilizing clear-sky and cloudy-sky Earth radiance measurements. These measurements fall within the same spectral range as the short-wave infrared (SWIR) and exhibit sensitivity to clouds.
%%%%%%%%%%%%%%%%%%%%%%%%%%%%%%%%%%%%%%%%%%%%%%%%%%%%%%%%%
\subsection{Formaldehyde} \label{tmas:ch2o} Formaldehyde is an intermediate gas in nearly all oxidation chains of non-methane volatile organic compounds (NMVOCs), ultimately contribute to CO$_2$ production. NMVOCs are, together with NO$_x$, CO and CH$_4$, among the most important precursors of tropospheric O$_3$. The main source of CH$_2$ in the atmosphere is CH$_4$ oxidation, mentioned to be an O$_3$ depleting substance in chapter \ref{litreview}. Over the continents, the oxidation of higher NMVOCs emitted from vegetation, fires, traffic and industrial sources results in important and localized enhancements of the HCHO levels. The seasonal and inter-annual variations in formaldehyde distribution are primarily linked to changes in temperature and fire events, as well as shifts in anthropogenic activities. HCHO concentrations in the boundary layer can be directly related to the release of short-lived hydrocarbons, which mostly cannot be observed directly from space.
%%%%%%%%%%%%%%%%%%%%%%%%%%%%%%%%%%%%%%%%%%%%%%%%%%%%%%%%%
\subsection{Clouds} \label{tmas:clouds} Cloud properties are derived from the OCRA and ROCINN algorithms, which are currently applied to GOME and GOME-2 products. OCRA measures cloud fraction using UV and visible-light spectral ranges, while ROCINN retrieves cloud height, pressure, and optical thickness (or albedo). Version 3.0 of the algorithms is employed, incorporating a more realistic approach that treats clouds as optically uniform layers of light-scattering particles. Furthermore, cloud parameters are provided for a cloud model that assumes the cloud behaves as a Lambertian reflecting surface. The resulting estimated cloud energy was created using transformations depicted in Chapter \ref{methods:thermo}.
%%%%%%%%%%%%%%%%%%%%%%%%%%%%%%%%%%%%%%%%%%%%%%%%%%%%%%%%%
\subsection{Tropospheric O$_3$} \label{tmas:S5Pozone} For this TROPOMI product, there are two algorithms used to deliver total column O$_3$ values: GDP for near-real-time and GODFIT for offline products. This thesis focuses on the use of the GDP algorithm, as it is currently being used to generate the operational total O$_3$ products of GOME, SCIAMACHY and GOME-2. Since 2017, these products have served as the final data source in most atmospheric chemistry-based machine learning models and have also informed policy-making for the EPA and EMWCF (\cite{schultz_tropospheric_2017}). This product serves as the baseline for O$_3$ values over the AOI. In contrast to earlier detections of O$_3$, this monitor reports values at $\frac{mol}{m^2}$ and does not require the transformation mentioned in Chapter \ref{methods:tomi}.
%%%%%%%%%%%%%%%%%%%%%%%%%%%%%%%%%%%%%%%%%%%%%%%%%%%%%%%%%%%%%%%%%%%%%%%%%%%%%%%%%%%%%%%%%%%%%%%%%%%%%%%%%%%%%%%%%%%%
\section{VIIRS Nighttime Lights} \label{tmas:ntl} There are two data sets for night time lights. This thesis used the Suomi National Polar Orbiting Partnership (SNPP) Visible Infrared Imaging Radiometer Suite (VIIRS) which supports a Day-Night Band (DNB) sensor that provides global daily measurements of nocturnal visible and near-infrared (NIR) light. These are suitable for Earth system sciences and applications. VNP46A2 is the short name for the daily moonlight- and atmosphere-corrected nighttime light (NTL) product called VIIRS/NPP Gap-Filled Lunar BRDF-Adjusted Nighttime Lights Daily L3 Global available at a 500 m \parencite{roman2018}. The VIIRS DNB's ultra-sensitivity in low light conditions enables the generation of high-quality night-time light imagery that utilizes substantial improvements to sensor resolution and calibration. 

Compared to the previous era of the Defense Meteorological Satellite Program/Operational Linescan System (DMSP/OLS), the updated VIIRS provides enhanced monitoring of nighttime phenomena and anthropogenic light emissions. This paper used the SNPP/VIIRS DNB dataset to characterize the surface dynamics of NO$_2$ columns that come from anthropogenic sources. This offered extremely minimal linear correlation, and NDVI was found to represent this process slightly better. Nighttime lights imagery from VIIRS is available at a daily resolution; data before 2012 exists at yearly intervals from 2000 to 2012 via 1 km resolution \parencite{zhao2020}. The Consistent and Corrected Nighttime Lights (CCNL) dataset reprocessed from the Defense Meteorological Program's (DMSP) Operational Line-Scan System (OLS) Version 4. Although inter-annual inconsistency, saturation, and blooming effects have improved the data quality of this historical dataset, further corrections are needed to provide daily representations of historical nighttime light data.
%%%%%%%%%%%%%%%%%%%%%%%%%%%%%%%%%%%%%%%%%%%%%%%%%%%%%%%%%%%%%%%%%%%%%%%%%%%%%%%%%%%%%%%%%%%%%%%%%%%%%%%%%%%%%%%%%%%%
\section{MODIS Vegetation Indices} \label{tmas:ndvi} The primary vegetation indices from the Moderate Resolution Imaging Spectroradiometer MODIS, generated at 16-day intervals and across multiple spatial resolutions, provide consistent spatial and temporal comparisons of vegetation canopy greenness. This composite property encompasses leaf area, chlorophyll content, and canopy structure \parencite{ramapriyan2022}. The primary vegetation index (VI) utilized in this analysis is the normalized difference vegetation index (NDVI). The Enhanced Vegetation Index (EVI) was gathered, but yielded less correlation to DAMO$_3$ than NDVI. NDVI can be constructed from Landsat imagery as well: \[NDVI = \frac{band_{R}-band_{NIR}}{band_{R}+band_{NIR}}\] where band$_{R}$ and band$_{NIR}$ represent the detections of red and near-infrared wavelengths, respectively.

These vegetation indices are derived from daily, atmospherically corrected, bidirectional surface reflectance. These indices ensure continuity with NOAA's AVHRR NDVI time series, thereby supporting historical and climate applications. This NDVI representation utilizes a MODIS-specific compositing method based on product quality assurance metrics to eliminate low-quality pixels. Since the MODIS sensors onboard the Terra and Aqua satellites are identical, the vegetation index algorithm can only generate composites at 16-day intervals, with composites produced 8 days apart. 
%%%%%%%%%%%%%%%%%%%%%%%%%%%%%%%%%%%%%%%%%%%%%%%%%%%%%%%%%%%%%%%%%%%%%%%%%%%%%%%%%%%%%%%%%%%%%%%%%%%%%%%%%%%%%%%%%%%%
\section{Final Surface O$_3$ Datasets} \label{tmas:findatasets} The correlation and normality of each potential feature were checked before modeling. Feature data missing between dates due to errors in satellite imagery were interpolated using the raster techniques in Chapter \ref{methods}. DAMO$_3$ values were interpolated for monitors that had less than 5 percent missing days over the initial time frame. The results of DAMO$_3$ interpolation can be seen in Figure \ref{tab:interpolation}; where linear regression, polynomial, and the Akima methods were tested. These missing values likely stem from either monitor malfunction or movement, as many historical monitors are still in use today. \begin{table}[h!]
    \centering
    \begin{tabular}{ccccc}
        \hline
        \textbf{Method} & \textbf{mse} & \textbf{mae} & \textbf{R$^2$} & \textbf{rmse} \\
        \hline
        Linear   &  0.03 & 4.16 &  0.79 & 5.49 \\
        ModAkima &  0.03 & 4.22 &  0.78 & 5.56 \\
        Akima 1D &  0.03 & 4.23 &  0.78 & 5.57 \\
        P(0)     &  0.04 & 4.67 &  0.74 & 6.04 \\
        P(2)     &  2.53 & 48.38 &-16.95 & 50.35 \\
        P(1)     &  2.55 & 48.62 &-17.05 & 50.48 \\
        P(3)     &  2.63 & 48.09 &-17.62 & 51.27 \\
       \hline
    \end{tabular}
    \caption{Results of DAMO$_3$ Interpolation.}
    \label{tab:interpolation}
\end{table} The interpolation methods resulted in an average error of around $\pm{4.19}$ ppb using a combination of linear and modified Akima \parencite{akima_new_1970}. These methods were well received, as \ref{fig:varsot} indicates that values over time were predicted within the minimum and maximum of the overall data set. Table \ref{tab:interpolation} presents the basic statistics of each interpolation strategy, wherein linear-based methods accounted for approximately 79 percent of DAMO$_3$. Polynomial methods were the worst when applied to DAMO$_3$; however, these methods are still implemented due to their high predictive power for features like NO$_2$ and NDVI. Although these values were not missing within this short time span, an old test script with models ranging from 2010 to 2024 necessitated the inclusion of these methods in the final script. 

The Pearson correlation coefficient values over time, as illustrated in Figure \ref{fig:histograms} and detailed in Table \ref{tab:variable_codes} were considered in the creation of each feature combination. Incorporating the feature changes noted in Chapter \ref{methods:featsfm} facilitated the integration of highly correlated variables into distinct CTM-based trends.  Although 44 features were extracted, only 31 were utilized in total, with a maximum of 24 employed simultaneously for each ensemble. The strongest correlations, as indicated in Figure \ref{fig:scatters} were used to create four unique datasets for this project. This allowed historical data (H.D) to have multiple high correlation variables similar to that of modern data (M.D) yet with satellite data dating back to the early 1990s \parencite{hersbach2020}. Each dataset was designed to incorporate features that were not selected for other datasets, with the exception of seasonal dummy variables and mean vapor pressure deficit. The latter, representing molecular water volume, was included in each dataset by happenstance.

Seasonal dummy variables were included to each data set in addition to their chosen features. Fall had the least correlative power among them. This was consistent with trends mentioned in Chapter \ref{subs:sofat1} and similar works (e.g., \cite{li_large_2021}; \cite{seroji_ground_2016}; \cite{zhao_assessing_2016}). Many of the initially created monthly variables were too redundant outside of the winter and summer seasons. Due to their weaker correlations, the monthly dummy variables were omitted from the final model. The Pearson correlation coefficients for the features included in each model illustrate how the resulting complex trend may manifest post-training.
%%%%%%%%%%%%%%%%%%%%%%%%%%%%%%%%%%%%%%%%%%%%%%%%%%%%%%%%%
\subsection{Historical Data} \label{tmas:hd} The most significant linear associations identified from the ERA5 and TOMS/OMI datasets were utilized due to their historical relevance. To establish the historical dataset (H.D) for this thesis, features that represented weekly moving averages and high linear associations were selected. The primary meteorological drivers of this dataset included vapor pressure deficit, burn index, NDVI, and specific humidity. Additionally, weekly moving averages of downward shortwave radiation and maximum surface temperature were incorporated. This dataset is likely to exhibit high multicollinearity among its features, as illustrated in Figure \ref{fig:histcorr}. 
\label{tmas:td}
\begin{figure}[h!]
    \centering
    \includegraphics[width=0.75\linewidth]{ThesisFilePkg//figs/hist_corrs.png}
    \caption{Pearson Correlation of Historical Data with DAMO$_3$}
    \label{fig:histcorr}
\end{figure} H.D. will primarily focus on the representation of meteorological conditions influencing DAMO$_3$, with NDVI being the sole connection to anthropogenic drivers. As a chemical connection, the KE estimate of the TOMS/OMI imagery data was included. This dataset also incorporates the least negatively correlated features, resulting in the lowest complexity for each model. The aim with this dataset was to develop the best model utilizing features available before 2000. High-quality NDVI imagery from MODIS, discussed in Chapter \ref{tmas:ndvi} is available at the specified 16-day intervals for this project, and via 
%%%%%%%%%%%%%%%%%%%%%%%%%%%%%%%%%%%%%%%%%%%%%%%%%%%%%%%%%
\subsection{Modern Data} \label{tmas:md} Modern data (M.D) consists of S5P imagery and chemical features that interact with DAMO $_3$. This data includes the KE estimate of S5P O$_3$, surface NO$_2$, water column energy, formaldehyde total column density (CH$_2$O), estimated surface CH$_2$O, carbon monoxide (CO), and the S5P estimated total column O$_3$ temperature. The correlations for M.D are illustrated in Figure \ref{fig:moderncorr}. 
\begin{figure}[h!]
    \centering
    \includegraphics[width=0.75\linewidth]{ThesisFilePkg//figs/modern_corrs.png}
    \caption{Pearson Correlation of Modern Data with DAMO$_3$}
    \label{fig:moderncorr}
\end{figure} Most of M.D features are derived from S5P imagery with the exception of mean vapor pressure deficit. This dataset also contained the same number of negatively correlated features as the larger dataset outlined in chapter \ref{tmas:goat24}, introducing additional complexity for each model compared to H.D. The intent with this dataset was to create the best model indicative of chemical interactions with anthropogenic sources. 
%%%%%%%%%%%%%%%%%%%%%%%%%%%%%%%%%%%%%%%%%%%%%%%%%%%%%%%%%
\subsection{Literature-Based Data} \label{tmas:td} The associations identified with O$_3$ in Chapter \ref{chap:litreview} were employed as the Literature-based Dataset (L.B.D), irrespective of their correlation with DAMO$_3$. These variables were primarily selected due to their temporal variability in relation to DAMO$_3$. The dataset includes surface NO$_2$, cloud fraction, NDVI, S5P kinetic energy, burn index and weekly moving averages of average N/E wind speed and S5P total column O$_3$ temperature. The correlations among these variables are depicted in Figure \ref{fig:theorycorr}.
\begin{figure}[h!]
    \centering
    \includegraphics[width=0.75\linewidth]{ThesisFilePkg//figs/theory_corrs.png}
    \caption{Pearson Correlation of Theoretical Data with DAMO$_3$}
    \label{fig:theorycorr}
\end{figure}
Many of the transformed variables were included due to their high correlations and representation of meteorological conditions with DAMO$_3$. This dataset also contained the most negatively correlated features, providing additional complexity for each model. The goal of this dataset was to inspect the incorporation of non-linear processes and their impacts on the final proposed method. While similar to M.D, the inclusion of NDVI in L.B.D makes it the most complex process in this thesis. The likelihood of strict trend establishment from these features is relatively low due to the equal amounts of positive and negatively correlated features.
%%%%%%%%%%%%%%%%%%%%%%%%%%%%%%%%%%%%%%%%%%%%%%%%%%%%%%%%%
\subsection{Best Correlating Features} \label{tmas:goat24} The top 24 features correlating with DAMO$_3$ were used at once to make the final predictive dataset. This was done to consider twice the number of features as the H.D, M.D, and T.D datasets with MLP, which tends to yield better performance with an extensive feature set. The outcomes of presenting correlation coefficients are illustrated in Figure \ref{fig:goat24corr}. The full dataset had the most diversity concerning DAMO$_3$ interactions. The selected features included vapor pressure deficit, surface NO$_2$, S5P total column O$_3$, estimated temperature, stratospheric NO$_2$, a weekly moving average of maximum surface temperature, minimum surface temperature, weekly moving average of mean pressure deficit, surface CH$_2$O, burn index, weekly moving average of average N/E wind speeds, water column energy, water column density, weekly moving average of the DU transformed DAMO$_3$ representation, cloud fraction, and CH$_2$O slant density. 
\begin{figure}[h!]
    \centering
    \includegraphics[width=0.75\linewidth]{ThesisFilePkg//figs/goat24_corrs.png}
    \caption{Pearson Correlation of the Top 24 Features with DAMO$_3$}
    \label{fig:goat24corr}
\end{figure} This dataset contained fewer negatively correlated features compared to T.D as NDVI, precipitation, and NTL exhibited low linear correlations. It comprises meteorological variables such as downward solar radiation, cloud fraction, and burn index, as well as chemical variables including surface NO $_2$, H$_2$O column number density, and KE estimates derived from both tropospheric O$_3$ collections. It utilized all Weekly moving averages including the maximum NE wind speed, DU conversion metric, and mean vapor pressure deficit. Most of the transformed variables were included due to their high correlation and representation of meteorological conditions with DAMO$_3$. This dataset exhibited slightly less complexity as NDVI was omitted. The goal was to develop a feature set of larger scale to analyze the impact on each estimated complex trend. 